% SPDX-License-Identifier: Apache-2.0
\section{Memory to Pixels, Across Two Clocks}
\label{sec:x-scanout}

\S\ref{sec:x-dma} reads memory in bursts and \S\ref{sec:x-cdc}
carries words from one clock to another. A video scanout is both at
once, and this is the two of them joined with a third part at the far
end:

\begin{center}\footnotesize
\begin{tabular}{@{}l@{}}
\code{Rom -- AxiPer -- AxiHost -- LineFetch ]--- ChanCdc ---[ LineBuf} \\
\hspace{1.1cm}\textit{bus clock, period four}
\hspace{1.3cm}\textit{pixel clock, period six} \\
\end{tabular}
\end{center}

\code{LineFetch} reads a line from memory on the bus clock,
\code{ChanCdc} carries the words to the pixel clock, and
\code{LineBuf} holds them for the raster to read a column at a time.
All three are parts in \code{//lib/parts} and all three are lowered.
The wiring in the example is the only thing here that is not
hardware.

\subsection{Why the line buffer registers its output}

\code{LineBuf} could read the memory straight out to the port, and
the first version did. It is wrong twice over.

A block RAM has an output register and a design that does not use it
is asking the memory to settle combinationally in the time a pixel
lasts. And a raster wants a pixel that is stable for the whole cycle
it is shown, not one that moves when the column moves. So
\code{shown} holds the pixel the column named at the last edge, and
the port is that register rather than the memory.

The cost is one cycle of lead, which a raster has: it knows which
column it will want next.

\lstinputlisting[rangebeginprefix=//\ begin\{,rangebeginsuffix=\},%
rangeendprefix=//\ end\{,rangeendsuffix=\},includerangemarker=false,%
linerange=linebuf-linebuf,caption={\code{lib/parts/src/dma.rs}: the
far end of the scanout.},label={lst:x-scanout}]{lib/parts/src/dma.rs}

\subsection{What the run checks}

Thirty-two pixels, read back a column at a time on the pixel clock,
each one the word that was at its own address in memory. The memory
holds its address in each word, so a pixel that came from the wrong
burst, crossed badly, or landed in the buffer out of step says which
place it came from rather than merely being wrong.

The two clocks are deliberately not in a simple ratio: period four
and period six with a phase of two. A crossing whose clocks tick
together is not crossing anything, and edges that stay in step are
how a crossing fault hides.

\subsection{The fault this example found}

The netlist could not be checked against the run when this was
written, and the reason was in the checking rather than in any of the
three parts.

A plain input was read from the trace one of its own clock's periods
late, so a column set between two pixel edges reached the entity an
edge early and \code{shown} moved one cycle before it should. Three
testbench visits make up one pixel cycle, and there were 31 column
changes after the first, which is the 93 differences the log
showed. That is issue 405, and it is a sibling of issue 384: 384 was
a handshake being compared where the trace holds nothing, this was a
level being applied before its edge.

Neither could appear under the default clock, where every tick the
testbench visits is an edge of every port's clock. They needed a unit
with ports on a slower one, which is what a scanout is.
